%% Theory of Computing
%% File: v022a007.tex
%% Publication date: October 11, 2026
%% Authors: Emanuele Viola

%
%
%
%

%
%
%
%
%
%
%
%
%
%
%
%


\documentclass{toc}

\usepackage{orcidlink}   %

\tocdetails{%
  volume=22, number=7, year=2026, firstpage=1,
  received={July 13, 2020},
  revised={January 13, 2021, January 29, 2026},
  published={October 11, 2026},        %
  doi = {10.4086/toc.2026.v022a007},   %
}

%
%

%
%
%
%
%
%
%
\newtheorem{thm}[theorem]{Theorem}
%
\newtheorem{rem}[theorem]{Remark}   %
%
\newtheorem{lem}[theorem]{Lemma}  %
%
%
\newtheorem{cor}[theorem]{Corollary}
\newtheorem{challenge}[theorem]{Challenge}
\newtheorem{defn}[theorem]{Definition}
%
%

\DeclareMathOperator{\sign}{\mathrm{sign}}  %
\DeclareMathOperator{\Maj}{\mathrm{Maj}}  %
\DeclareMathOperator{\MAJ}{\mathrm{MAJ}}  %
\DeclareMathOperator{\Ltf}{\mathrm{Ltf}}  %
\DeclareMathOperator{\Time}{\mathrm{Time}}  %
\DeclareMathOperator{\NTime}{\mathrm{NTime}}  %
\DeclareMathOperator{\AC}{\mathrm{AC}}  %
\DeclareMathOperator{\ACC}{\mathrm{ACC}}  %
\DeclareMathOperator{\NP}{\mathrm{NP}}  %
\DeclareMathOperator{\PPP}{\mathrm{P}}  %
\DeclareMathOperator{\EEE}{\mathrm{E}}  %
\DeclareMathOperator{\CKT}{\mathrm{CKT}}  %

\newcommand{\zo}{\{0, 1\}}
\newcommand{\F}{\mathbb{F}}
\renewcommand{\P}{\mathbb{P}}
\newcommand{\E}{\mathbb{E}}
\newcommand{\poly}{\mathrm{poly}}
\newcommand{\polylog}{\mathrm{polylog}}
\newcommand{\eps}{\epsilon}
\renewcommand{\a}{\alpha}
\renewcommand{\b}{\beta}
\newcommand{\e}{\epsilon}

%
\runningauthor{Emanuele Viola}
%
\runningtitle{New Lower Bounds for Probabilistic Degree}

\copyrightauthor{Emanuele Viola}

\begin{document}
\begin{frontmatter}

\title{New Lower Bounds for Probabilistic Degree and AC0 with Parity
  Gates\titlefootnote{A prelimiary version of this paper appeared
    in \href{https://eccc.weizmann.ac.il/report/2020/015/}{ECCC as TR20-015
    (2021)} \cite{eccc-version}.}}
%

%
\tocpdftitle{New Lower Bounds for Probabilistic Degree and AC0 with Parity Gates}


\tocpdfauthor{Emanuele Viola}

%
%
\author[viola]{Emanuele Viola\thanks{Supported by NSF CCF award 1813930.}
\,\,\raisebox{-0.17ex}{\orcidlink{0000-0001-6091-1824}}}
%

\begin{abstract}

We make the first progress on probabilistic-degree lower bounds and
correlation bounds for polynomials since the papers by Razborov and
Smolensky in the 80's. The bounds hold for computing some function
$f:\zo^{n}\to\zo$ in $\EEE^{\NP}$, and include:

(1) $\Omega(n/\log^{2}n)$ lower bounds probabilistic degree. This
is optimal up to a factor $O(\log^{2}n)$. The previous best lower
bound was $\Omega(\sqrt{n})$ proved in the 80's by Razborov and Smolensky.

(2) $\exp(\Omega(n/\log^{2}n)^{1/(h-1)})$ lower bounds on the size
of depth-$h$ $\AC^{0}[\oplus]$ circuits, for any $h$. This almost
matches the $\exp(\Omega(n^{1/(h-1)}))$ lower bounds for $\AC^{0}$ by
%
H{\aa}stad.  %
The previous best lower bound was $\exp(\Omega(n^{1/(h+1)}))$
by Rajgopal, Santhanam, and Srinivasan
(MFCS'18)  %
who recently improved Razborov
and Smolensky's $\exp(\Omega(n^{1/(2h-2)}))$ bound.

(3) $(1/2-(\log^{O(h)}s)/n)$ average-case hardness for size-$s$
depth-$h$ $\AC^{0}[\oplus]$ circuits under the uniform distribution.
The previous best was $(1/2-(\log^{O(h)}s)/\sqrt{n})$. A concurrent
%
paper   %
by Chen and Ren
(STOC'20)  %
obtains an incomparable result.

The %
previous best lower bounds in (1) and (3),
mentioned above,  %
held for the
Majority function. Each of the new lower bounds in this paper is false
for Majority. For (2) the previous best held for $\EEE^{\NP}$.

The proofs build on %
%
%
Williams's
``guess-and-SAT'' method. For (1)
we show how to use a PCP by Ben-Sasson and Viola 
(ICALP'14)   %
towards probabilistic-degree
lower bounds. For (3) we combine a recent %
method %
by Alman and Chen
(STOC'19, SICOMP 2025)  %
with hardness amplification.
\end{abstract}

\tockeywords{lower bound, probabilistic degree, AC0, PCP, polynomials}

%
\tocacm{F.2.3}    %

%
\tocams{68Q17}    %

\end{frontmatter}

\section{Introduction}
%
%
%
%

%
%
%
\subsection{History, motivation}  %
\label{preamble}                  %

\emph{Probabilistic degree} is a fundamental complexity measure of
boolean functions that has been intensely studied since it was introduced
by Razborov~\cite{Raz87}. Probabilistic degree measures how well
a function can be computed by a ``random'' polynomial, from a suitable
distribution. One can consider polynomials over different fields.
For simplicity in this paper we focus on the field $\F_{2}$ with
two elements, and we only mention here that the results can be extended
to $\F_{p}$ for prime~$p$.

\begin{defn}
The \emph{$\e$-error degree} of a function $f:\zo^{n}\to\zo$ is
the minimum $d$ such that there is a distribution $D$ on polynomials
over $\F_{2}$ of degree $\le d$ such that
$\P[D(x) \ne f(x)] \le \eps$
for every $x\in\zo^{n}$.
\end{defn}

The parameter $\e$ ranges from $0$ to at most $1/2$, because outputting
a random bit yields error $1/2$. In this %
article  %
we will consider both
the setting when $\e$ is close to $0$, and the setting when $\e$
approaches $1/2$. The interplay between the two settings will be
important.

%
%
Razborov and Smolensky \cite{Raz87,Smo87,Smolensky93} proved
that the Majority function on $n$ bits has constant-error degree
$\Omega(\sqrt{n})$.%
%
%
\footnote{In this paper, %
the notation $O$ and $\Omega$ %
implies positive %
absolute constants.}
Better lower bounds have not been obtained despite much
research. One difficulty is that many classes of functions do have
$O(\sqrt{n})$ probabilistic degree.
Alman and Williams
%
show \cite[Theorem 1.2]{DBLP:conf/focs/AlmanW15}   %
that any symmetric function on $n$ bits has $\e$-error
polynomials of degree $O(\sqrt{n\log(1/\e)})$. Earlier, Srinivasan
\cite{DBLP:conf/fsttcs/Srinivasan13} proved a slightly weaker bound.
The recent %
article~\cite{srinivasan2019probabilistic}  %
gives a nearly
tight characterization of the probabilistic degree of symmetric functions.
Srinivasan %
\cite[Theorem 12]{DBLP:conf/fsttcs/Srinivasan13}  %
also
proved that threshold functions (with arbitrary weights) have $\e$-error
polynomials of degree 
$\sqrt{n}(\log n\cdot\log1/\e)^{O(1)}$.
Golovnev, Kulikov, and Williams~%
%
\cite{DBLP:journals/eccc/GolovnevKW18}
conjecture that DeMorgan formulas of size $s$ have probabilistic
degree $O(\sqrt{s})$. %
%
%

\medskip

Razborov and Smolensky actually proved a tradeoff between the error
and the degree.
%
%
%
%
For some constant $c$ and any degree $\le c \sqrt{n}$, %
%
%
%
%
%
%
%
we can get constant error as mentioned above,
but for lower degree we can push the error
to be close to $1/2$. Formally, the tradeoff gives that for any
degree $d$ the error is
\begin{equation}
%
%
%
\e\ge1/2- O(d/\sqrt{n})\,.\label{eq:smolensky}   %
\end{equation}
%
%
%
%

This tradeoff is tight for Majority %
\cite{Viola-L12requiresCor},  %
and is
the best available for an explicit function, for any degree $d\ge\log_{2}n$.
In particular, conceivably every explicit function (say in $\PPP$ or
even $\NP$) has %
$(1/2-1/\sqrt{n})$-error %
degree $\log_{2}n$.
%
For small degrees, $d<\log_{2}n$, the iterated Cauchy--Schwarz argument
of~\cite{BNS92} gives stronger error bounds.   %
%
%
%
%
%
%
%
See \cite{ViolaGF2,ViW-GF2,viola-FTTCS09}
for direct arguments.  %
%
However in this paper we are only concerned with $d\ge\log_{2}n$.

We note that by Yao's duality principle~\cite{DBLP:conf/focs/Yao77},
the claim that a function $h$ has $\e$-error degree $d$ is equivalent
to the claim that there exists a distribution over the inputs such
that any fixed degree-$d$ polynomial fails to compute $h$ with probability
$\e$ over inputs drawn from that distribution. In the special case
where the latter distribution is uniform we call the function $h$
$\e$-\emph{hard }for polynomials of degree $d$.
\begin{defn}
\label{def:hardness}Let $h:\zo^{n}\to\zo$ be a function, and let
$F$ be a set of boolean functions on $n$ bits. We say that $h$
is $\e$-hard for $F$ if for every $f\in F$ we have $\P_{x}[f(x)\ne h(x)]\ge\e$,
where $x$ is uniform in $\zo^{n}$.
\end{defn}

We remark that the state of lower bounds is the same for probabilistic
degree and hardness against %
polynomials; no  %
better lower bounds are
known if we allow a non-uniform distribution. Hardness results for
polynomials are surveyed in %
\cite[Chapter 1]{viola-FTTCS09}.   %

\medskip

Probabilistic degree is linked to a wide range of other problems.
In particular, progress on probabilistic degree is \emph{necessary}
for progress on circuit lower bounds, rigidity, and communication
complexity. For some of these problems, it is also \emph{sufficient}.
We now elaborate on this. We begin by observing that, with regards
to lower bounds, improving the
%
%
%
\expeqref{tradeoff}{eq:smolensky} %
for small $d$ is easier than improving it for large $d$. Specifically,
if a function $f$ has $(1/2-\e)$-error degree $d$ then it has constant-error
degree $O(d/\e)$ (see \expref{Lemma}{lem:connection-tradeoff}).

%
\subsubsection{Circuit lower bounds for \texorpdfstring{$\AC^{0}[\oplus]$}{AC0(+)}}

The circuits $\AC^{0}[\oplus]$ consist of And, Or, and Parity gates
with unbounded fan-in, and have constant depth. This class lies at
the frontier of the techniques in circuit lower bounds. The
%
%
work  %
by Razborov and Smolensky used probabilistic-degree lower bounds to
obtain lower bounds of $\exp(\Omega(n^{1/(2h-2)}))$ on the size of
$\AC^{0}[\oplus]$ circuits of depth $h$. This bound was recently
improved to $\exp(\Omega(n^{1/(h+1)}))$ by Rajgopal, Santhanam, and
Srinivasan~\cite{DBLP:conf/mfcs/RajgopalS018}, see 
also~\cite{DBLP:conf/coco/OliveiraS019}.
Still, these lower bounds are lower than the $\exp(\Omega(n^{1/(h-1)}))$
that can be obtained for $\AC^{0}$ via
%
%
H{\aa}stad's switching lemma~\cite{Has87}.  %

The gap between $\AC^{0}$ and $\AC^{0}[\oplus]$ lower bounds is closely
related to the state of probabilistic-degree lower bounds: If we could
prove linear lower bounds on constant-error degree, that is eliminate
the square root in 
\expeqref{Equation}{eq:smolensky}, %
then we would obtain lower
bounds for $\AC^{0}[\oplus]$ that are as strong as the $\AC^{0}$ ones,
using the same reasoning
as  %
in the proof of \expref{Corollary}{cor:size}.

Also, we don't have strong average-case lower bounds. It is consistent
with our knowledge that for every explicit function there is a polynomial-size
$\AC^{0}[\oplus]$ circuit computing it correctly on $1/2+1/\sqrt{n}$
fraction of the inputs $\zo^{n}$. Note that degree-$d$ polynomials
are a special case of depth-2 $\AC^{0}[\oplus]$ circuits of size $n^{O(d)}$,
so this is similar (but not equivalent to) the fact that we don't
have such strong average-case lower bounds for degree $\log n$.

\subsubsection{Rigidity}
\label{sec:rigidity}   %

%
Nearly half a century ago  %
Valiant \cite{Val77} asked to construct
matrices that cannot be approximated by matrices %
of low rank.  %
We can say that a $\zo^{n}\times\zo^{n}$ matrix $M$ is $\e$-error
rigid for rank $r$ if for any rank-$r$ matrix $F$ we have 
$\P_{i,j}[M_{i,j}\ne F_{i,j}]\ge\e$,
where $i$ and $j$ are uniform. Viewing matrices as truth-tables
of functions, this is the same as  %
being  %
$\e$-hard for rank-$r$ functions.
Again this can be considered over different fields and we focus on
$\F_{2}$ in this %
article.  %
The rigidity question is motivated by a host
of applications ranging from circuit lower bounds to communication
complexity. Non-explicitly, there exist matrices which are constant-error
rigid for $r=\Omega(2^{n})$. But despite intense research the available
constructions are far from this. Recently, an exciting %
paper  %
%
%
%
by Alman and Chen~\cite{AlmanC19} 
constructed in the complexity class $\Time(2^{O(n)})^{\NP}=\EEE^{\NP}$
matrices $M$ which are constant-error rigid for rank $2^{n^{1/4-\e}}$,
for any $\e>0$. They also gave a conditional construction for higher
rank $2^{n^{1-\alpha}}$ for all $\alpha>0$. For a discussion of
previous rigidity bounds and applications we refer the reader to~\cite{AlmanC19}
and Lokam's survey~\cite{Lokam09}.

Probabilistic-degree lower bounds stand in the way of further progress.
Indeed, as pointed out by Servedio and Viola~\cite{ServedioV-sparse},
sparse polynomials are a special case of low-rank matrices. Specifically,
write the truth-table of a polynomial on $2n$ variables as a $2^{n}\times2^{n}$
matrix (partitioning the variables in half arbitrarily). Now observe
that each monomial is a rank one matrix. Hence a polynomial with $s$
monomials yields a rank $s$ matrix. Because a polynomial of degree
$\sqrt{2n}$ has $\le n^{O(\sqrt{n})}$ monomials, we see that better
probabilistic-degree lower bounds are necessary for improving the
results in~\cite{AlmanC19} to any rank $2^{\omega(\sqrt{n}\log n)}$.
(We can remove the $\log$ factor when considering \emph{sparsity,
}see~\cite{ServedioV-sparse} for discussion. Here we focus on \emph{degree}
for simplicity.)

\subsubsection{Communication complexity}

In the influential number-on-forehead communication model introduced
by Chandra, Furst, and Lipton~\cite{CFL83}, $k$ parties collaboratively
wish to compute a function on $n$ bits. The bits are divided in $k$
blocks (figuratively corresponding to the foreheads) and the twist
is that Party $i$ knows all blocks except the %
$i$-th.   %
For this model,
we have no lower bounds when $k\ge\log n$ (for $k<\log n$ they were
first proved in~\cite{BNS92}, see \cite{ChT93,Raz00,ViW-GF2} for
expositions). %
%
%
%
%
%
For background on communication complexity see the books
\cite{KuN97,RaoY2019}.

Interestingly, it turns out that proving communication lower bounds
for more parties requires better probabilistic-degree lower bounds.
We include in this paper a stronger quantitative version of this connection
from~\cite{Viola-map}, see \expref{Theorem}{thm:NOF-connection}, which
already kicks in for $k=O(\log^{2}n)$ parties exchanging $O(\log^{3}n)$
bits of communication. In particular, improving the \expeqref{tradeoff}{eq:smolensky}
is necessary to rule out such protocols.

\subsubsection{Other circuit classes}

Many papers in the literature study $\AC^{0}[\oplus]$ circuits augmented
with other types of gates. In fact, average-case lower bounds for
$\AC^{0}[\oplus]$ are closely related to worst-case lower bounds for
the circuit class $\Maj\circ \AC^{0}[\oplus]$ of such circuits with
a majority gate at the output~\cite{HMPST93}. Correspondingly, we
only have lower bounds when the fan-in of the majority gate is $t<n/\poly\log(n)$.
Such a lower bound follows simply by approximating the entire circuit
by a polynomial of degree $\sqrt{n}/\poly\log(n)$, and then using
the probabilistic degree lower bound. For the approximation, we proceed
by approximating Majority using~\cite{DBLP:conf/focs/AlmanW15} and
each $\AC^{0}[\oplus]$ circuit using~\cite{Raz87}, and then compose
the polynomials.

Alman and Chen~\cite{AlmanC19} consider the class of 
$\AC^{0}[\oplus]\circ \Ltf_{t}\circ \AC^{0}[\oplus]\circ \Ltf$
circuits, where $\Ltf$ are threshold functions with arbitrary weights,
and prove that the fan-in $t$ of the middle $\Ltf$ gates satisfies
$t\ge n/\log^{O(h)}s$ for depth-$h$ and size-$s$ circuits that
compute some function in $\EEE^{\NP}$. A lower bound for this class could
not just rely on probabilistic degree, because the $\Ltf$ gates at the
input already require $\sqrt{n}$ probabilistic degree.

Related classes of circuits were studied by Alman, Chan, and Williams
($\ACC^{0}\circ \Ltf{}_{n^{2-\Omega(1)}}\circ \Ltf$,
%
\cite[Corollary 1.1]{AlmanCW16}),   %
and Tamaki (depth-2 circuits with $n^{2}/\poly\log(n)$
gates, where each gate can be any symmetric function or an $\Ltf$~\cite{Tamaki16}).

\subsubsection{Pseudorandom generators}

Another important motivation for improving the \expeqref{tradeoff}{eq:smolensky}
and average-case lower bounds for $\AC^{0}[\oplus]$ is the construction
of \emph{pseudorandom generators.} Nisan's landmark paper~\cite{Nis91}
showed how to use a $(1/2-\e)$-hard function to obtain approximately
$1/\e$ bits of pseudorandomness. Hence, a suitable improvement to
the tradeoff -- for the specific case of uniform distributions --
would yield better generators for low-degree polynomials \cite{BoV-gen,Lov09,Viola-d},
a long-standing problem for which several new approaches have recently
been proposed \cite{DBLP:conf/coco/ChattopadhyayHH18,DBLP:conf/innovations/ChattopadhyayHL19,DBLP:journals/corr/abs-2008-01316,Viola-L12requiresCor}.
Nisan's approach yields pseudorandom generators with small error,
in which case improving probabilistic degree lower bounds is known
to be necessary~\cite{Viola-d}. The situation with constant-error
generators is less clear, cf.~\cite{viola-PRG-requires-MAJ}. We
also remark that other pseudorandom-generator constructions that have
been proposed do not go through for classes such as $\AC^{0}[\oplus]$,
see \cite{ShV-dec,GrinbergSV-adaptivemajority,viola-PRG-requires-MAJ}.

\medskip

We consider it worthwhile to try to identify an open problem which
is as simple as possible while at the same time being a bottleneck
for as wide a range of complexity lower bounds as possible. The above
discussion suggests that probabilistic degree may be a good candidate
for this.

\subsection{Our results}

In this %
paper  %
we show that $\EEE^{\NP}$ requires %
nearly linear %
constant-error
degree, improving on the $\Omega(\sqrt{n})$ lower bounds established
by Razborov and Smolensky \cite{Raz87,Smo87,Smolensky93} for the
majority function.
\begin{thm}
\label{thm:prob-degree-lower-bound}There is a function $f:\zo^{n}\to\zo$
in $\EEE^{\NP}$ with $1/3$-error degree $\Omega(n/\log^{2}n)$, for
infinitely many $n$.
\end{thm}

The bound in \expref{Theorem}{thm:prob-degree-lower-bound} is tight up
to a factor of $O(\log^{2}n)$ because every function on $n$ bits
has polynomials of degree $n$ with error $0$. We note that, as mentioned
earlier, the bound in \expref{Theorem}{thm:prob-degree-lower-bound} is
false for any symmetric or threshold function
\cite{DBLP:conf/fsttcs/Srinivasan13,DBLP:conf/focs/AlmanW15}.

\expref{Theorem}{thm:prob-degree-lower-bound} has several corollaries
stated next. First, as mentioned earlier, it actually improves the
\expeqref{tradeoff}{eq:smolensky} for every setting of parameters (and
not just constant error).
\begin{cor}
\label{cor:tradeoff}The function in \expref{Theorem}{thm:prob-degree-lower-bound}
has $(1/2-\e)$-error degree $\Omega(\e n/\log^{2}n)$, for infinitely
many $n$ and any $\e$.
\end{cor}

So in particular we can set $\e=(\log^{4}n)/n$ and still get degree
%
$\Omega(\log^{2}n)$ while  %
from the \expeqref{tradeoff}{eq:smolensky}
we could not set $1/2-\e$ larger than $1/2-1/\sqrt{n}$. Error larger
than $1/2-1/\sqrt{n}$ was not even available for polynomials of degree
$\log n$.

%

Then we improve the size lower bounds for $\AC^{0}[\oplus]$.
\begin{cor}          %
\label{cor:size}Any depth-$h$ $\AC^{0}[\oplus]$ circuit computing
$f$ in \expref{Theorem}{thm:prob-degree-lower-bound} has size
%
%
\[ \ge\exp(\Omega(n/\log^{2}n)^{1/(h-1)})\,, \]  %
for infinitely many $n$.
\end{cor}

This improves on the $\exp(\Omega(n^{1/(h+1)}))$ lower bounds
%
%
by Rajgopal, Santhanam, and Srinivasan~\cite{DBLP:conf/mfcs/RajgopalS018}.
The bound in \expref{Corollary}{cor:size} is false for any symmetric
function~\cite{DBLP:conf/coco/OliveiraS019}, and it almost matches
the $\exp(\Omega(n^{1/(h-1)}))$ lower bounds that are available for
$\AC^{0}$ via switching lemmas~\cite{Has87}.

Next we prove a general lower bound for $\AC^{0}[\oplus]$ circuits
suitably augmented with gates computing threshold or arbitrary symmetric
functions. We show that we can allow such gates as long as the product
of their fan-in's along any root-leaf path in the circuit is at most
%
$n^{2}/\poly\log(n)$.   %
First we formally define the class.
\begin{defn}
The class $t$-SoT-$\AC^{0}[\oplus]$ consists of circuits made of
And, Or, Parity, Sym, and $\Ltf$ gates such that the product of the fan-in
of the Sym and $\Ltf$ gates along any root-leaf path in the circuit is
at most $t$. Here each Sym gate computes an arbitrary symmetric function,
and each $\Ltf$ gate computes a threshold function with arbitrary weights.
\end{defn}

\begin{cor}
\label{cor:SOT-lower-bounds} Any $t$-SoT-$\AC^{0}[\oplus]$ circuit
of depth $h$ and size $s$ computing the function $f$ in \expref{Theorem}{thm:prob-degree-lower-bound} requires $t\ge n^{2}/\log^{O(h)}s$,
for infinitely many $n$.
\end{cor}

In particular, $\Maj_{t}\circ \AC^{0}[\oplus]$ circuits of size $s$
and depth $h$ have $t\ge n^{2}/\log^{O(h)}s$, improving on the previous
best of $t\ge n/\log^{O(h)}s$ which follows from Razborov~\cite{Raz87}
and Smolensky \cite{Smo87,Smolensky93} and the probabilistic degree
of Majority \cite{DBLP:conf/fsttcs/Srinivasan13,DBLP:conf/focs/AlmanW15}.

Also, $\AC^{0}[\oplus]\circ \Ltf{}_{t}\circ \AC^{0}[\oplus]\circ \Ltf$
circuits of size $s$ and depth $h$ have $t\ge n/\log^{O(h)}s$.
The previous best was $t\ge\sqrt{n}/\log^{O(h)}s$ recently proved
by Alman and Chen~\cite{AlmanC19}. Our proof of this result simply
approximates the circuit and then invokes the probabilistic degree
lower bound. By contrast Alman and Chen rely on a combination of degree
and rank arguments.

\subsubsection{Uniform distribution}

\expref{Corollary}{cor:tradeoff} implies that for any distribution $C$
on polynomial-size constant-depth $\AC^{0}[\oplus]$ circuits one has
$\P_{C}[C(x)\ne f(x)]\ge1/2-\poly\log(n)/n$ for some fixed input
$x$. This is because as mentioned earlier the functions computable
by such circuits of %
polynomial size  %
have poly-logarithmic probabilistic
degree~\cite{Raz87}. As mentioned earlier, Yao's duality principle
\cite{DBLP:conf/focs/Yao77} then guarantees hardness under \emph{some
}distribution, not guaranteed to be uniform. With a different argument,
we show how to strengthen this to hold under the uniform distribution.
\begin{thm}
  \label{thm:uniform} Let $h$ be an integer and $s$ a
monotone increasing  %
function satisfying
%
%
$s(\Omega(n\log n)) < 2^{O(\log^{2}s(n))}$  %
%
%
%
%
(such as $s(n)=2^{n^{\alpha}}$
for any
positive  %
$\alpha$).%
\footnote{When an inequality involves multiple occurrences of the $O$ and $\Omega$
  asymptotic notation, we interpret each such symbol as a positive constant,
  existentially quantified.  For instance, for positive-valued functions $f,g$ and
  monotone increasing positive-valued functions $s,t$, the inequality
  $s(\Omega(f(n)))\le t(O(g(n)))$ means $(\exists c_1, c_2 > 0)
  (s(c_1 f(n)) \le t(c_2 g(n))$ for all sufficiently large $n$.}
%
%
%
%
Then there %
is a function $f:\zo^{n}\to\zo$ in $\EEE^{\NP}$ that is $(1/2-(\log^{O(h)}s)/n)$-hard
for size-$s$ depth-$h$ $\AC^{0}[\oplus]$ circuits, for infinitely
many $n$.
\end{thm}
%

Again, the previous best was hardness $1/2-(\log^{O(h)}s)/\sqrt{n}$,
under any distribution. \expref{Theorem}{thm:uniform} gets us closer
to a correlation bound that could be used to obtain new pseudorandom
generators via Nisan's method~\cite{Nis91}.

\subsection{Techniques \label{subsec:Techniques}}

The results in this paper rely on a masterful method by Williams~\cite{Williams13-Improving}
which we call the ``guess-and-SAT'' method. This method has been
used to prove several lower bounds for computing functions in $\EEE^{\NP}$, NEXP, NQP,
and related classes; lower bounds which we do not know how prove by
other means. Researchers have established lower bounds against several
non-uniform circuit classes starting with
%
Williams's  %
$\ACC^{0}$ lower
bound~\cite{Williams14-acc-jacm}, see \cite{Williams11,Williams13,Williams14,AlmanCW16,Tamaki16,DBLP:journals/eccc/Chen19,DBLP:conf/latin/ChenOS18,DBLP:conf/stoc/MurrayW18,DBLP:conf/mfcs/RajgopalS018,AlmanC19,VyasWilliams20,ChenRen20}.

To explain the proof of \expref{Theorem}{thm:prob-degree-lower-bound}
let us quickly review the method. Suppose we aim to prove a lower
bound for computing a function in $\EEE^{\NP}$ by a circuit class $\CKT$.
We begin by considering a language $L$ in non-deterministic time
$2^{n}$ which cannot be solved in non-deterministic time $o(2^{n})$.
Such a language exists by the non-deterministic time-hierarchy theorem.
We note that this hierarchy only holds for infinitely many input lengths,
as opposed to all sufficiently large input lengths, which is why all
these lower bounds, including the ones in this paper, only hold for
infinitely many input lengths.

We reduce $L$ to 3SAT with about $2^{n}$ clauses (and variables)
using the Cook-Levin theorem. This is done in a specific way -- using
\emph{indexing circuits} -- which is clarified below. We then consider
the function $f$ that on input $x$ and $i$ computes the $i$ bit
of a canonical satisfying assignment corresponding to $x$, if it
is the case that $x\in L$. One can show $f\in \EEE^{\NP}$. Towards a
contradiction we assume that $f$ has a small circuit $C\in \CKT$.
The algorithm for $L$ \emph{guesses }$C$ and then combines it with
the \emph{indexing circuit} that given $z$ of about $n$ bits outputs
the variables in clause $z$ of the 3CNF (and with a post-process
circuit which possibly negates the outcomes, and outputs Or). This
gives a new circuit $C'$ such that it suffices to know if $C'$ accepts
every input $z$ to decide membership in $L$. If this can be done
in time $o(2^{n})$, we contradict the hierarchy theorem.

Actually, the original proof of the $\ACC^{0}$ lower bound does not
work this way, due to the lack of reductions to 3SAT with $\ACC^{0}$
indexing circuits. This modular version of the method has been obtained
in~\cite{JahanjouMV-explicit}.

\subsubsection{Using PCP}

In the original paper~\cite{Williams13-Improving} it was already
suggested that rather than reducing to 3SAT we can reduce to gap-3SAT,
or in other words use \emph{probabilistically checkable proofs} (PCP).
In a PCP we refer to the input $z$ of about $n$ bits as the \emph{randomness}.

Again, existing PCPs did not come equipped with an indexing circuit
that was efficient enough for a modular version of the method as above.
This was remedied by Ben-Sasson and Viola in~\cite{Ben-SassonV-AC0PCP}
building on \cite{Ben-SassonGHSV05,Ben-SassonS08}. Their PCP is critical
for this work. In this PCP, the indexing circuit simply computes \emph{projections}
(a.k.a. 1-local functions). In other words, on input the randomness
$z$, the PCP queries the proof at locations $q_{1},q_{2},\ldots,q_{t}$
where each $q_{i}$ is about $n$ bits, and each bit of $q_{i}$ depends
on only one bit of $z$.

Projections work well for probabilistic degree, because they do not
increase degree. However, the PCP makes a polynomial number $t$ of
queries which are then post-processed by a 3CNF $\phi$. It is not
known if a similar PCP with only a constant number of queries exists,
see discussion in~\cite{Ben-SassonV-AC0PCP}. This 3CNF $\phi$ on
a polynomial number of variables is problematic, because it may not
have small degree. (Note that we cannot approximate it by a polynomial,
as we could get a bad approximation which accepts wrong proofs, and
we would have no way of detecting it.) Another issue is that we are
trying to prove a lower bound against a \emph{distribution }on polynomials.
Typical instantiations of the method work against \emph{fixed} circuits.

\subsubsection{Our approach}

To go around these issues we begin by modifying slightly the PCP in
\cite{Ben-SassonV-AC0PCP}. Rather than making a polynomial number
of queries and applying the 3CNF $\phi$ to their answers, the modification
simply picks a random clause $i$ in $\phi$, queries its three variables,
negates the outputs as needed, and outputs Or.

This PCP will still accept correct proofs with probability $1$, because
the 3CNF evaluates to $1$ and so all clauses are true.

If on the other hand the proof is incorrect, by the soundness of the
PCP~\cite{Ben-SassonV-AC0PCP} $\phi$ evaluates to $0$ on at least
say a $0.5$ fraction of the choices for the randomness $z$. For
each such $z$, at least one of the clauses is not satisfied. Because
the number of clauses is a polynomial, we have at least a polynomial
probability, over the choice of $i$, to pick a false clause. Hence,
over the choice of both $z$ and $i$, this PCP will accept an incorrect
proof with probability $\le1-0.5/n^{b}$ for some fixed $b$.

The gap between the accept and reject probabilities is small. However,
it is going to be large enough for us. Recall that, towards a contradiction,
we are assuming that the proof of the PCP is going to be computed
by a low-degree polynomial. In fact, by a distribution on polynomials
which errs with constant probability. Now an important point is that
we are allowed to \emph{reduce this error}. This is possible because
we have a distribution on polynomials that has an advantage on every
fixed input. So we can just sample many times from the distribution
and compute majority to drive the error down exponentially. Specifically,
by sampling $c\log n$ many times, for a constant $c$ depending on
$b$, we can drive the error to $\e<0.5/n^{b}$.

At this point by an averaging argument it is possible to fix a polynomial
$p$ that maintains a gap between the accept and reject probabilities
of the PCP. Hence, for inputs in the language $L$, there exists a
proof given by a polynomial $p$ such that the PCP will still accept
with probability $>1-0.5/n^{b}$, while for inputs not in $L$, the
PCP will accept any proof with probability $<1-0.5/n^{b}$.

The algorithm for $L$ proceeds by guessing the polynomial $p$ and
combining it with the PCP. So if we can give an algorithm running
in time $o(2^{n})$ to compute the acceptance probability of this
PCP exactly, we derive the desired contradiction.

This acceptance probability cannot yet be written as the number of
satisfying assignments to a low-degree polynomial. This is because
the map from $i$ to the clause is not known to be of low degree.
However, there are only a polynomial number of choices for $i$. So
we simply enumerate over all of them. For a fixed $i$, we can write
the acceptance probability of the PCP as the number of satisfying
assignments to a low-degree polynomial in $z$.

\medskip

Summarizing, we have reduced proving probabilistic-degree lower bounds
to computing the number of satisfying assignments to a polynomial.

\medskip

To finish the proof of \expref{Theorem}{thm:prob-degree-lower-bound},
we need to compute the number of satisfying assignments to a degree-$d$
polynomial in $n$ variables in time $2^{n-\Omega(n/d)+O(\log n)}$.
Such an algorithm was recently obtained by Williams: it is a special
case of Theorem 26 in~\cite{DBLP:conf/coco/Williams18}.~\cite{DBLP:conf/soda/LokshtanovPTWY17}
obtain running time $2^{n-\Omega(n/d)+o(n)}$ which is not quite sufficient
for our purposes. However the results in these papers are more general
and so the proofs are somewhat involved. We give a streamlined presentation
of the algorithm that we need, using a fairly straightforward combination
of techniques that have been used in the corresponding algorithms
for $\ACC^{0}$, \emph{Orthogonal Vectors, }and \emph{systems of polynomial
equations} \cite{Williams14-acc-jacm,DBLP:conf/soda/ChanW16,DBLP:conf/soda/LokshtanovPTWY17}.
This concludes the overview of the proof of \expref{Theorem}{thm:prob-degree-lower-bound}.

The corollaries then follow using known simulations of circuit classes
by probabilistic polynomials.

\subsubsection{Proof of \expref{Theorem}{thm:uniform}}

To prove the lower bound under the uniform distribution
in \expref{Theorem}{thm:uniform} we use a different argument. Here our starting
point is a \emph{conditional }result by Alman and Chen~\cite{AlmanC19}.
Under the assumption that non-deterministic quasi-polynomial time
has polynomial-size circuits (\ie, $NQP\subseteq P/\poly$) they
show how to construct $2^{n}\times2^{n}$ matrices which are constant-error
rigid for rank $2^{n^{1-\alpha}}$ for any $\alpha>0$. Recall that
(the functions whose truth tables are) such matrices are also $\Omega(1)$-hard
for polynomials of degree $n^{1-\alpha-o(1)}$.

In this paper we also give a streamlined exposition of their result
for the case of polynomials, achieving slightly better parameters
than what claimed in~\cite{AlmanC19} (see \expref{Section}{sec:Proof-of-Lemma-stronger-AC}).

At the high level, \expref{Theorem}{thm:uniform} is established by combining
this result by Alman and Chen with \emph{hardness amplification},
a technique to transform hard functions into harder functions. For
background we refer the reader to Chapter 17 ``Hardness Amplification
and Error Correcting Codes'' in the textbook~\cite{SanjeevBarak2007},
and to the discussion in~\cite{ShV-dec}.

However, the use of hardness amplification is not straightforward.
One issue is that hardness amplification in general cannot be used
for classes such as $\AC^{0}[\oplus]$, see~\cite{GrinbergSV-adaptivemajority}
and the discussion in \cite{ShV-dec,Viola-PhDthesis-journal-hack}.
Our proof will use hardness amplification \emph{twice}, for rather
different purposes. Let us give an overview. Towards a contradiction
let us assume that every function in $\EEE^{\NP}$ has $\AC^{0}[\oplus]$
circuits of size $n^{c}$ which compute the function correctly on
a $1/2+1/n^{0.9}$ fraction of the inputs. By worst-case to average-case
hardness amplification %
\cite{BFNW93,Imp95,STV01}, %
%
this means
that $\EEE^{\NP}$ has small \emph{circuits, }not necessarily of constant
depth. However, this is sufficient to satisfy the assumption in the
result by Alman and Chen, which then yields a function which is $\Omega(1)$-hard
for polynomials of degree $n^{1-\alpha-o(1)}$, for any $\alpha$.
In particular, this function is $\Omega(1)$-hard for $\MAJ_{n^{2-2\alpha-o(1)}}\circ \AC^{0}[\oplus]$
circuits of size $n^{c'}$. An important point is that here we can
set $c'>c$, which is necessary for the final contradiction. Now in
our second application of hardness amplification we use a fine analysis
of Yao's famous XOR lemma~\cite{GNW95} which follows from Impagliazzo's
theory of hard-core sets~\cite{Imp95} and essentially appears in
a beautiful paper by Klivans~\cite{Kli01}. The analysis essentially
says that if a function $f$ is $\Omega(1)$-hard for $\MAJ_{t^{2}}\circ C$
circuits then $XOR$ of $O(\log t)$ independent copies of $f$ is
$(1/2-1/t)$-hard for circuits $C$. Applying this analysis we produce
another function in $\EEE^{\NP}$ which violates our initial assumption
and concludes the proof.

We note that these uses of hardness amplification change the input
length, and so we require some mild conditions on the size function
$s=s(n)$ to finish the argument. (One can actually relax the conditions
at the price of having a more complicated statement.)

\subsubsection{Comparison with rigidity lower bound in~\cite{AlmanC19}}

We mentioned earlier that Alman and Chen %
\cite{AlmanC19} gave   %
an unconditional construction of $2^{n}\times2^{n}$ matrices which are
constant-error rigid for rank $2^{n^{1/4-\e}}$. This result does
not imply new probabilistic-degree lower bounds. For the latter, we
would need to obtain rank larger than $2^{\sqrt{n}}$.

It is natural to ask whether the approach in~\cite{AlmanC19} can
yield new probabilistic-degree lower bounds. We now explain what the
problem is and why rank $2^{\sqrt{n}}$ is in fact a natural barrier
for their approach.~\cite{AlmanC19}'s argument uses PCPs twice. The
first application is similar to what we described above. The second
application, following~\cite{DBLP:conf/coco/ChenW19}, is a PCP of
proximity that helps in computing the acceptance probability of the
first PCP.

The critical point is this: when using this approach to prove a lower
bound against a set of functions (for example polynomials of degree
$d$) \emph{this second PCP is applied to evaluate a circuit which
contains, and so is at least as big as, one of these functions}. This
makes the proof length of this PCP at least as long as the description
of the functions. In the case of polynomials of degree $d$, this
would be at least $n^{d}\ge2^{d}$. In particular, the randomness
used by the second verifier has length $u$ at least $d$. Now, the
algorithm in~\cite{AlmanC19} involves enumerating over all $2^{u}$
choices for this randomness, and running for each choice an algorithm
that counts the number of satisfying assignments. For the latter problem,
there are algorithms running in time $2^{n-\Omega(n/d)}$ 
(see \expref{Section}{subsec:=000023SAT-poly}). The overall
%
%
%
%
running time  %
accounting for
going over the $2^{u}$ choices and running the counting algorithm
on each is then $2^{d}\cdot2^{n-\Omega(n/d)}.$ However, this is only
better than $2^{n}$ when $d\le\sqrt{n}$.

The paper~\cite{AlmanC19} also contains an argument which they call
bootstrapping\emph{.} As discussed in~\cite{AlmanC19}, bootstrapping
runs into difficulties when the rank is larger than $\sqrt{n}$; the
same difficulties arise for degree larger than $\sqrt{n}.$

\subsubsection{Comparison with the work by Rajgopal, Santhanam, and 
Srinivasan~\cite{DBLP:conf/mfcs/RajgopalS018}}

As mentioned earlier, Rajgopal, Santhanam, and Srinivasan~\cite{DBLP:conf/mfcs/RajgopalS018}
recently proved $\exp(\Omega(n^{1/(h+1)}))$ lower bounds against
depth-$h$ $\AC^{0}[\oplus]$ circuits. Their proof does not rely on
the PCP from~\cite{Ben-SassonV-AC0PCP}, instead it uses a reduction
to 3SAT from the same paper which refines~\cite{JahanjouMV-explicit}
(the latter paper suffices if we replace $h+1$ with $h+2$ in the
bound). The use of~\cite{Ben-SassonV-AC0PCP} is straightforward while
the bulk of the proof in~\cite{DBLP:conf/mfcs/RajgopalS018} lies
in obtaining new satisfiability algorithms for $\AC^{0}[\oplus]$ circuits.
It is interesting to note that the balance in our proof 
of \expref{Corollary}{cor:size} is opposite. We rely on the PCP 
from~\cite{Ben-SassonV-AC0PCP}
to establish a new reduction from probabilistic-degree lower bounds
to computing the number of satisfying assignments to a polynomial,
while the algorithm for the latter task is fairly straightforward.

\subsection{Towards probabilistic rank}

As mentioned earlier, probabilistic degree is a special case of rigidity.
The notion of rigidity is in fact the same as that of hardness (\expref{Definition}{def:hardness}) for low-rank matrices. Razborov introduced~\cite{Razborov-rigidity}
a variant of rigidity which is to rigidity what probabilistic degree
is to hardness for low-degree polynomials: we have a distribution
on low-rank matrices that on every entry has the correct value except
with probability $\e$. By Yao's duality principle, this is the same
as correlating with respect to \emph{every }distribution. Following
the literature we call this notion \emph{probabilistic rank}. Probabilistic
rank has been studied by several works including \cite{DBLP:journals/cc/Wunderlich12,AlmanW17}.
This notion is particularly interesting because lower bounds for probabilistic
rank, besides being a prerequisite for rigidity lower bounds, also
suffice for some of the applications of rigidity lower bounds. For
example they suffice for communication lower bounds, see~\cite{DBLP:journals/cc/Wunderlich12},
and for $\AC^{0}[\oplus]\circ \Ltf_{t}\circ \AC^{0}[\oplus]\circ \Ltf$
lower bounds, see \cite{AlmanC19,AlmanW17}. For an application to
circuits operating on matrices see Jukna's book~\cite{Jukna2012},
Chapter 12.8. With a simple proof, in \expref{Section}{sec:Data-structures}
we also show that sufficiently strong lower bounds on probabilistic
rank imply new \emph{data-structure} lower\emph{ }bounds\emph{,} a
connection which as far as we know is new.

It is natural to ask if \expref{Theorem}{thm:prob-degree-lower-bound}
can be extended to bound probabilistic rank as well. From our approach
it follows that \emph{either} a strengthening of the PCP~\cite{Ben-SassonV-AC0PCP}
\emph{or }a strengthening of the algorithms for counting orthogonal
vectors \cite{DBLP:conf/soda/AbboudWY15,DBLP:conf/soda/ChanW16} would
give improved lower bounds for probabilistic rank.

On the PCP side we would need a PCP that on input randomness $z$
makes queries of the type $z\oplus a_{i}$ where $a_{i}$ are fixed
strings, and $\oplus$ is bit-wise XOR. One of the components of the
PCP~\cite{Ben-SassonV-AC0PCP} does have this form, but another (the
PCP of proximity for Reed Solomon codes, which remains the most intricate
part of the proof) does permute bits.

On the side of orthogonal vectors,
%
%
it would suffice to solve the following problem.
\begin{problem}
\label{prob:Given-six-functions}Given six functions $f_{1}',f_{1}'',f_{2}',f_{2}'',f_{3}',f_{3}'':\zo^{n}\to\zo^{r}$
and six subsets $s_{1}',s_{1}'',s_{2}',s_{2}'',s_{3}',s_{3}''\subseteq\{1,2,\ldots,2n\}$
of size $n$ each, compute
\[
\P_{z\in\zo^{2n}}\left[\sum_{i=1}^{3}\langle f'_{i}(z|_{s'_{i}}),f''_{i}(z|_{s''_{i}})\rangle=1\mod2\right]
\]
in time $2^{2n-\Omega(n/\log r)}$ when $r\le2^{\alpha n}$ for a
small enough $\alpha>0$. Here $z|_{s}$ denotes the $n$ bits of
$z$ indexed by $s$.
\end{problem}

The connection to rank is that a rank-$r$ $2^{n}\times2^{n}$ matrix
can be written as the truth-table of a function $f:\zo^{n}\times\zo^{n}\to\zo^{r}$
defined as $g(x,y):=\langle f'(x),f''(y)\rangle$ where $\langle,\rangle$
is inner product. Thus \expref{Problem}{prob:Given-six-functions} considers
the sum of three low-rank functions of \emph{different projections
of $2n$ bits. }As pointed out in~\cite{AlmanC19}, when the projections
are the same this can be reduced to a single matrix, and for that
one can apply algorithms for \emph{orthogonal vectors}~\cite{DBLP:conf/soda/ChanW16}.
It is not clear what happens for different projections, which could
arise in the PCP~\cite{Ben-SassonV-AC0PCP}.

Here's a special case that we can't yet solve. Divide the $2n$ bits
$z$ into 4 blocks $z=(x,y,v,w)$ of $n/2$ bits each. Can one compute
\[
\P_{z\in\zo^{2n}}[\langle f(x,y),f(v,w)\rangle+\langle f(x,v),f(y,w)\rangle+\langle f(x,w),f(y,v)\rangle=1\mod2]
\]
in time $2^{2n-\Omega(n/\log r)}$?

\subsection{Concurrent and subsequent work}

\subsubsection{Concurrent work}

A concurrent paper by Chen and Ren~\cite{ChenRen20} obtains a result
incomparable to \expref{Theorem}{thm:uniform}. They obtain functions
in $\EEE^{\NP}$ that are $(1/2-1/f(n))$-hard for $\AC^{0}[\oplus]$ (or
even $\ACC^{0}$) circuits of size $f(n)$. They implicitly prove this
for up to \emph{sub-half-exponential }functions $f(n)$, and the techniques
in~\cite{ChenRen20} do not apply to circuits of larger size. In particular,
for exponential-size circuits~\cite{ChenRen20} do not obtain better
hardness than the classic $1/2-1/\sqrt{n}$ result. Our techniques
apply even to exponential-size circuits, though our improvement on
hardness is smaller. For example, \expref{Theorem}{thm:uniform} yields
hardness $1/2-1/n^{0.99}$ for circuits of size $2^{n^{\Omega(1)}}$.
This state of affairs raises the exciting possibility of making both
the hardness and the size parameter exponential.

\subsubsection{Subsequent work}

The paper~\cite{ChenLW2020-almost} simultaneously improves the parameters
in~\cite{ChenRen20} and in \expref{Theorem}{thm:uniform} (improving
the function $f(n)$ above to exponential). It also shows how to make
the lower bounds ``almost everywhere'' instead of ``infinitely
often.'' The paper~\cite{DBLP:journals/corr/abs-2005-03123} constructs
rectangular (and smooth) PCPs and uses them to show the existence
of rigid matrices in E$^{\NP}$. The PCP we asked about above are in
particular rectangular. The results in this paper are obtained using
as black-box the PCP~\cite{Ben-SassonV-AC0PCP}. Both~\cite{DBLP:journals/corr/abs-2005-03123}
and~\cite{Ben-SassonV-AC0PCP} in turn rely on PCPs in~\cite{Ben-SassonGHSV05}.
Quantitatively,~\cite{DBLP:journals/corr/abs-2005-03123} prove hardness
$1/2-\Omega(1)$ against functions whose corresponding $2^{n}\times2^{n}$
matrix has rank $2^{n/O(\log n)}$. Recall
%
(see for example the argument in \expref{Section}{sec:rigidity}) %
that degree-$d$ polynomials have rank $\le O(\binom{n}{d})\le O((\eee n/d)^{d})$.
Hence~\cite{DBLP:journals/corr/abs-2005-03123} implies constant-error
probabilistic-degree lower bounds for degree up to $\Omega(n/(\log n\cdot\log\log n))$,
improving the degree bound in our main \expref{Theorem}{thm:prob-degree-lower-bound}
by almost a $\log n$ factor.

Building on both~\cite{ChenLW2020-almost} and~\cite{DBLP:journals/corr/abs-2005-03123},
the paper~\cite{Huang-Viola-average-case-rigidity} proves strong
average-case hardness $1/2-2^{-\Omega(k)}$ against functions whose
corresponding $2^{n}\times2^{n}$ matrix has rank $2^{n/O(k(k+\log n))}$
(for $k\le c\sqrt{n}$). This generalizes both~\cite{ChenLW2020-almost}
and~\cite{DBLP:journals/corr/abs-2005-03123}. The latter can be recovered
setting $k=O(1)$.

\subsubsection{Organization}

In \expref{Section}{sec:Probabilistic-degree-lower} we prove the probabilistic-degree
lower bound \expref{Theorem}{thm:prob-degree-lower-bound}. The corollaries
are then proved in \expref{Section}{sec:Proofs-of-corollaries}. The hardness
result under the uniform distribution, \expref{Theorem}{thm:uniform}
is proved in \expref{Section}{sec:Uniform-distribution:-Proof}. This
latter theorem relies on a slightly stronger version of a result by
Alman and Chen specialized to polynomials. We give a streamlined exposition
of this in \expref{Section}{sec:Proof-of-Lemma-stronger-AC}. In \expref{Section}{sec:Number-on-forehead} we include a quantitative version of
a connection from~\cite{Viola-map} between communication protocols
and probabilistic degree. In \expref{section}{sec:Data-structures} we
describe the connection to data-structure lower bounds.

\section{Probabilistic degree lower bound \label{sec:Probabilistic-degree-lower}}

In this section we prove \expref{Theorem}{thm:prob-degree-lower-bound}.
We begin with two subsections, then prove the main theorem.

\subsection{PCP tools \label{subsec:PCP-tools}}

We need a simple variant of a PCP by Ben-Sasson and the author~\cite{Ben-SassonV-AC0PCP}.
This PCP in turn is a relatively minor modification of the PCP~\cite{Ben-SassonGHSV05}
in which the verifier's query indexes are just projections of the
verifier's randomness bits. In fact, for the application in this paper
it is sufficient to have soundness polynomially bounded away from
$1$, so one can dispense with the soundness amplification step in
\cite{Ben-SassonV-AC0PCP}.

First we state the result from~\cite{Ben-SassonV-AC0PCP}, then state
and prove our variant. Here ``algorithm'' refers to any standard
model of computation, such as random-access multi-tape Turing machines,
cf.~\cite{Williams14-acc-jacm}.
\begin{lem}[{\cite{Ben-SassonV-AC0PCP}}] \label{lem:BV-PCP}Let $M$ be an algorithm
running in time $T=T(n)\geq n$ on inputs of the form $(x,y)$ where
$|x|=n$. Given $x\in\{0,1\}^{n}$ one can output in time $\poly(n,\log T)$
circuits $Q:\zo^{r}\to(\zo^{r})^{t}$ for $t=\poly(r)$ and $R:\zo^{t}\to\zo$
such that:

\textbf{Proof length.} $2^{r}\le T\cdot\poly\log T$,

\textbf{Completeness.} If there exists $y$ such that $M(x,y)$ accepts
then there exists a map $\pi:\zo^{r}\to\zo$ such that for any $z\in\zo^{r}$
we have $R(\pi(q_{1}),\ldots,\pi(q_{t}))=1$ where $(q_{1},\ldots,q_{t})=Q(z)$,

\textbf{Soundness.} If no $y$ causes $M(x,y)$ to accept, then for
every map $\pi:\zo^{r}\to\zo$, at most $1/n^{10}$ fraction of the
$z\in\zo^{r}$ have $R(\pi(q_{1}),\ldots,\pi(q_{t}))=1$ where $(q_{1},\ldots,q_{t})=Q(z)$,

\textbf{Complexity.} $Q$ is a projection (a.k.a.~1-local), \ie,
each output bit of $Q_{i}$ is one input bit, the negation of an input
bit, or a constant; $R$ is a 3CNF.
\end{lem}

Next is our variant.
\begin{lem}
\label{lem:NEW-PCP} Let $M$ be an algorithm running in time $T=T(n)\geq n$
on inputs of the form $(x,y)$ where $|x|=n$. Given $x\in\{0,1\}^{n}$
one can output in time $\poly(n,\log T)$ a collection of $\poly(r)$
circuits $Q_{i}:\zo^{r}\to(\zo^{r})^{3}$ and $R_{i}:\zo^{3}\to\zo$
such that:

\textbf{Proof length.} $2^{r}\le T\cdot\poly\log T$,

\textbf{Completeness.} If there exists $y$ such that $M(x,y)$ accepts
then there exists a map $\pi:\zo^{r}\to\zo$ such that for any $z\in\zo^{r}$
and any $i\le\poly(r)$ we have $R_{i}(\pi(q_{1}),\pi(q_{2}),\pi(q_{3}))=1$
where $(q_{1},q_{2},q_{3})=Q_{i}(z)$,

\textbf{Soundness.} If no $y$ causes $M(x,y)$ to accept, then for
every map $\pi:\zo^{r}\to\zo$, at most a $1-1/r^{O(1)}$ fraction
of the pairs $(z,i)\in\zo^{r}\times[\poly(r)]$ have $R_{i}(\pi(q_{1}),\pi(q_{2}),\pi(q_{3}))=1$
where $(q_{1},q_{2},q_{3})=Q_{i}(z)$,

\textbf{Complexity.} Each $Q_{i}$ is a projection (a.k.a.~1-local),
\ie, each output bit of $Q$ is one input bit, the negation of an
input bit, or a constant; each $R_{i}$ is an Or of three literals.
\end{lem}

\begin{proof}
We start with the PCP in \expref{Lemma}{lem:BV-PCP}. Rather than computing
all of the 3CNF $R$, the verifier only checks that clause $i$ is
satisfied. The function $R_{i}$ is therefore an Or of three literals.
Completeness is immediate. To argue soundness, note that for every
choice of $z$ such that $R(\pi(q_{1}),\ldots,\pi(q_{t}))=0$, there
is at least one clause of $R$ that is not satisfied. Because the
number of clauses is $\poly(r)$, and by the soundness of the PCP
in \expref{Lemma}{lem:BV-PCP}, we have that the probability over $(z,i)$
that $R_{i}(\pi(q_{1}),\pi(q_{2}),\pi(q_{3}))=0$ is at least $(1-1/n^{10})\cdot(1/r^{O(1)})\ge1/r^{O(1)}$.
\end{proof}

\subsection{Counting roots of polynomials\label{subsec:=000023SAT-poly}}

In this subsection we show how to count efficiently the satisfying
assignments (a.k.a.~roots) of a low-degree polynomial. We give a
streamlined exposition of an algorithm which is a special case of
Theorem 26 in~\cite{DBLP:conf/coco/Williams18}.
\begin{thm}[{\cite[Theorem 26]{DBLP:conf/coco/Williams18}}]
\label{thm:=000023SAT-poly}
Given an $\F_{2}$-polynomial $p$ of degree $d$ in $n$ variables
we can compute the number of satisfying assignments (a.k.a.~non-roots)
in time $2^{n-\Omega(n/d)+O(\log n)}$.
\end{thm}

For the proof we use \emph{modulus-amplifying polynomials} from the
90's. We use them in a way similar to their use for orthogonal-vector
algorithms (\cite{DBLP:conf/soda/ChanW16} Theorem 1.2). However using
the latter algorithms directly appears to give a slightly worse bound.
(For example, the dimension of the corresponding vectors is $\Omega(n^{d})$
and so the running time would be $2^{n-\Omega(n/d\log n)}$ according
to~\cite{DBLP:conf/soda/ChanW16} Theorem 1.2.) Instead we use a classic
dynamic-programming approach going back to Yates in the %
1930's,  %
see \cite{Williams14-acc-jacm}.
\begin{lem}[{Modulus-amplifying polynomials, \cite{Yao90,BeT94}}] \label{lem:=00005BModulus-amplifying-polynomials=00005D}
For every integer $\ell$ there is a univariate polynomial $F_{\ell}$
of degree $2\ell-1$ over the integers such that for every input $y$:

- $y\equiv0\mod2\Rightarrow F_{\ell}(y)\equiv0\mod2^{\ell},$ and

- $y\equiv1\mod2\Rightarrow F_{\ell}(y)\equiv1\mod2^{\ell}.$

Without loss of generality, the coefficients of $F_{\ell}$ are in
$\{0,1,\ldots,2^{\ell}-1\}$. The coefficients of the polynomials
can be computed in time $\poly(\ell)$.
\end{lem}

%
\begin{lem}[{Yates in the 1930's, see~\cite{Williams14-acc-jacm}}] %
%
\label{lem:Yates}
Let $f:\zo^{n}\to\{0,1,\ldots,2^{\poly(n)}\}$. Given the truth-table
of $f$ we can compute in time $2^{n}\poly(n)$ the truth-table of
the function $g:\zo^{n}\to\mathbb{Z}$ defined as $g(x)=\sum_{y\subseteq x}f(y)$,
where we identify $\zo^{n}$ with the subsets of $\{1,2,\ldots,n\}$.
\end{lem}

A proof of this lemma via dynamic programming can be found in~\cite{Williams14-acc-jacm}.
\begin{proof}[Proof of \expref{Theorem}{thm:=000023SAT-poly}] Let $\ell:=\alpha n/d$
for an $\alpha>0$ to be determined later. View $p$ as a polynomial
over the integers modulo $2^{\ell+1}$ (with coefficients in $\zo$).
Let $F_{\ell+1}$ be the modulus amplifying polynomial of degree $2(\ell+1)-1$.
The idea is to compose $F_{\ell+1}$ and $p$ to obtain polynomial
$p'$ of degree $O(d\ell)=O(\alpha n)$. Because we are only interested
in the values $\zo$ for the variables, we can simplify $p'$ replacing
terms $x_{i}^{j}$ with $j\ge1$ with $x_{i}$. With this simplification,
the number of monomials in $p'$ will be $M:=\sum_{i=0}^{O(\alpha n)}\binom{n}{i}\le2^{H(O(\alpha))n}$
where $H$ is the binary entropy function. The coefficients of $p'$
have $O(\ell)$ bits.

Note that we can write down $p'$ in time $\poly(\ell)\cdot\poly(M)$.
For example we can perform the $\poly(\ell)$ operations corresponding
to $F_{\ell}$ to the polynomial $p'$ working modulo $2^{\ell+1}$
. After each operation, we simplify $p'$ as mentioned above.

Now define $p''(x_{1},x_{2},\ldots,x_{n-\ell}):=\sum_{x_{n-\ell+1},x_{n-\ell+2},\ldots,x_{n}}p'(x_{1},x_{2},\ldots,x_{n})$.

We can write down $p''$ in time $2^{\ell}\cdot\poly(M)\cdot\poly(\ell)$.
Note $p''$ has again $\le M$ monomials, and coefficients of $O(\ell)$
bits.

Note that for any input $x\in\zo^{n-\ell}$, $p''(x)$ counts the
number of satisfying assignments among the $2^{\ell}$ completions
of $x$ to an input of length $n$. This is because $p'$ has been
amplified to work modulo $2^{\ell+1},$ so there is no ``wrap-around''
when summing over $2^{\ell}$ inputs. Hence we only need to evaluate
$p''$ on all inputs. At this point we forget that $p''$ has low
degree and we just consider $p''$ as a function $f:\zo^{n-\ell}\to\{0,1,\ldots,2^{O(\ell)}\}$
mapping monomials to coefficients. More in detail, we consider $p''$
as a polynomial with exactly $2^{n-\ell}$ monomials, with possibly
zero coefficient. We view $p''$ as the function $f$ whose input
is interpreted as a monomial and whose output is the corresponding
coefficient. Evaluating $p''$ on every input $x$ is the same as
computing the truth-table of the function $g(x)=\sum_{y\subseteq x}f(y)$.
\expref{Lemma}{lem:Yates} solves the evaluation problem in time $2^{n-\ell}\cdot\poly(n)$.

Overall, the algorithm's running time is 
\[
2^{\ell}\poly(M)\cdot\poly(\ell)+2^{n-\ell}\cdot\poly(n)\le2^{O(H(O(\alpha)))n}+2^{n-\alpha n/d+O(\log n)}\,.
\]
For a small enough $\alpha$ the first term is less than the second,
and the result follows.
\end{proof}

\subsection{Proof of \expref{Theorem}{thm:prob-degree-lower-bound}}

Assume towards a contradiction that every function in $\EEE^{\NP}$ has
$1/3$-error polynomials of degree $d=\alpha n/\log^{2}n$, where
$\a$ is a constant to be determined later. First, we reduce the error.
Consider the majority of $t$ independent polynomials. This is a probabilistic
polynomial of degree $t\cdot d$, simply writing the majority exactly
as a degree-$t$ polynomial and composing polynomials. On every input,
by a Chernoff bound the error of this new polynomial is $2^{-\Omega(t)}$.
Taking $t=O(c\log n)$ we obtain that every $f\in \EEE^{\NP}$ on $n$
bits has $1/n^{c}$-error polynomials of degree $d':=O(dc\log n)$.
We will set $c$ later.

Let $L\in\mathrm{\NTime}(2^{n})\setminus\mathrm{\NTime}(o(2^{n}))$
\cite{Cook73,SeiferasFM78,Zak83}. We want to show that $L\in\mathrm{\NTime}(o(2^{n}))$
to reach a contradiction. Consider the algorithm $f(x,y)$ that on
input $x\in\zo^{n}$ and $y\in\{1,2,\ldots,2^{n}\poly(n)\}$ constructs
the circuits from the PCP in \expref{Lemma}{lem:NEW-PCP} corresponding
to the verifier for $L$, computes the first satisfying assignment
if one exists (otherwise computes say the all zero string), and outputs
its bit $y$. This algorithm can be implemented in $\EEE^{\NP}$ by computing
the satisfying assignment one bit at the time.

Recalling $(q_{1},q_{2},q_{3})=Q_{i}(z)$, where $z\in\zo^{r}$ and
$r=n+O(\log n)$, we have for every $x\in L$:
\[
\P_{z\in\zo^{r},i\in[r^{O(1)}]}[R_{i}(f(x,q_{1}),f(x,q_{2}),f(x,q_{3}))=1]=1;
\]

while for every $x\not\in L$ and every function $\pi$

\[
\P_{z,i}[R_{i}(\pi(q_{1}),\pi(q_{2}),\pi(q_{3}))=1]\le1-1/n^{b}
\]

for a constant $b=O(1)$.

By the argument above, $f$ has a $1/n^{c}$-error polynomial of degree
$d'$. We pick $c=b+1$. For a fixed input $x\in L$ we can hard-wire
the input $x$ in the probabilistic polynomials for $f$ and write
$P=P_{x}$ for the corresponding distribution of polynomials. Because
$P$ has error $1/n^{c}$ on every fixed input, and because each constraint
is an Or of three literals, we have for $x\in L$:
\[
\P_{z,i,P}[R_{i}(P(q_{1}),P(q_{2}),P(q_{3}))=1]\ge1-1/n^{c}\,.
\]

Note that the weaker bound $1-3/n^{c}$ follows by a union bound.
The stronger bound holds because it suffices that one input to $R_{i}$
is $1$ to have $R_{i}=1$. Also we can take each occurrence of $P$
inside the probability to be the same sample from the distribution
of polynomials. Neither observation is crucial here but may be useful
elsewhere.

By an averaging argument, there exists a fixed polynomial $p$ of
degree $d$ such that
\[
\P_{z,i}[R_{i}(p(q_{1}),p(q_{2}),p(q_{3}))=1]\ge1-1/n^{c}\,.
\]

On the other hand for any $x\not\in L$ and for every proof $\pi$
we have:
\[
\P_{z,i}[R_{i}(\pi(q_{1}),\pi(q_{2}),\pi(q_{3}))=1]\le1-1/n^{b}<1-1/n^{c}\,.
\]

In particular this holds if we set $\pi$ to be any polynomial $p$.

Hence, there is a gap between the probabilities. We now give an algorithm
to detect the gap. Specifically, we give an efficient non-deterministic
algorithm that decides if there exists a polynomial $p$ such that
$\P_{z,i}[R_{i}(p(q_{1}),p(q_{2}),p(q_{3}))=1]\ge1-1/n^{c}.$

On input $x$, the algorithm guesses a polynomial $p'$ of degree
$d'$ on $n'=n+O(\log n)$ variables. The number of bits required
to specify such a polynomial is
$\ell=\sum_{i=0}^{d'}\binom{n'}{i}\le(\eee n'/d')^{d'}\le n^{d'}=n^{O(dc\log n)}=2^{O(\alpha cn)}$
which is $o(2^{n})$ for a small enough $\alpha$ depending only on
$c$. The algorithm then computes the corresponding probability, and
accepts if it is at least $1-1/n^{c}$.

\emph{Computing the probability}. The algorithm will compute the probability
for every fixed value of $i$, and then average the values. Because
the number of $i$ is only $\poly(r)=\poly\log T=\poly(n)$, this
only incurs a polynomial slow-down.

Once $i$ is fixed, the queries $q_{1},q_{2},q_{3}$ are linear functions
of the input $z$. Composing these functions with $p'$ and writing
the output $R_{i}$ exactly as a degree-3 polynomial, we obtain a
polynomial $q$ of degree $O(d')$. To write down the polynomial we
need first to compute the query circuits $Q_{i}$ and the circuit
$R_{i}$. This can be done in time $\poly(n)$ by \expref{Lemma}{lem:NEW-PCP}.
Composing these with $p'$ can be done in time polynomial in the length
of the representation of $p'$ which is $\poly(\ell)$. As above,
the latter quantity can be made say $\le2^{n/2}$, again for $\alpha$
small enough depending on $c$. We then compute the number of inputs
that cause $q$ to accept using \expref{Lemma}{thm:=000023SAT-poly}.
The degree of $q$ is $O(d')\le O(\alpha cn/\log n)$. The running
time of the algorithm in \expref{Lemma}{thm:=000023SAT-poly} is $2^{n-\Omega(\log n)/(\alpha c)+O(\log n)}$.
For a small enough $\alpha$ the overall running time of the algorithm
is $o(2^{n})$, concluding the proof.

\section{Proofs of corollaries \label{sec:Proofs-of-corollaries}}

In this section we prove \expref{Corollaries}{cor:tradeoff}, \ref{cor:size},
\expref{and}{cor:SOT-lower-bounds}.  %
We need several upper bounds on probabilistic
degree from the literature.

\subsection{Probabilistic degree upper bounds}

We rely on Razborov's seminal approximation result~\cite{Raz87}.
We use the recent tighter version by
Kopparty and Srinivasan~\cite{DBLP:journals/toc/Kopparty018},
who also show their bound is nearly tight.
\begin{lem}[{\cite{DBLP:journals/toc/Kopparty018}}]
 \label{lem:approx-KS}Any $\AC^{0}[\oplus]$
circuit of depth $d$ and size $s$ has $\e$-error degree
%
%
at most
\[ O(\log s)^{d-1}\log(1/\e)\,. \]  %
\end{lem}

We also need upper bounds on the probabilistic degree of symmetric
and threshold functions, by Alman and Williams, and Srinivasan.
\begin{lem}[{\cite[Theorem 1.2]{DBLP:conf/focs/AlmanW15}}]\label{lem:AW-maj}
Any symmetric function on $n$ bits has $\e$-error
%
%
%
degree
at most  %
%
$O(\sqrt{n\log(1/\e)})$.
\end{lem}


\begin{lem}[{\cite[Theorem 12]{DBLP:conf/fsttcs/Srinivasan13}}]
\label{lem:Sri-LTF}
Any threshold function (a.k.a.~Ltf)
has $\e$-error degree $\le\sqrt{n}(\log n\cdot\log1/\e)^{O(1)}$.
\end{lem}


\subsection{Proofs of corollaries}

First we prove \expref{Corollary}{cor:tradeoff} by connecting two settings
of parameters for probabilistic degree, showing that improving the
state-of-the-art for small degree and error approaching $1/2$ is
easier than improving it for large degree and constant error. Then
we prove \expref{Corollary}{cor:SOT-lower-bounds}.
\begin{lem}
\label{lem:connection-tradeoff}Suppose that $f:\zo^{n}\to\zo$ has
$(1/2-\e)$-error degree $d$. Then $f$ has $1/10$-error degree
$O(d/\e)$.
\end{lem}

In particular, for $\e=1/\sqrt{n}$ and $d=\log n$ we obtain degree
$O(\log n\cdot\sqrt{n})$.
\begin{proof}
Let $P$ be a distribution witnessing $(1/2-\e)$-error degree $d$.
We take the majority of $t=O(1/\e^{2})$ independent copies of $P$.
By a Chernoff bound this majority has error $\le1/20$. Moreover,
Majority on $t$ bits has $1/20$-error polynomials of degree $O(\sqrt{t})$
\cite{DBLP:conf/focs/AlmanW15}. So replacing the majority with a
probabilistic polynomial gives error $1/20+1/20\le1/10$ and degree
$O(d/\e)$.
\end{proof}
From this the proof of \expref{Corollary}{cor:tradeoff} is immediate.
Now we prove the other two corollaries.
\begin{proof}[Proof of \expref{Corollary}{cor:size}]
	We apply \expref{Lemma}{lem:approx-KS}
with $\e=1/3$ to obtain that the function computed by the circuit
has $1/3$-error degree $O(\log s)^{h-1}$. By \expref{Theorem}{thm:prob-degree-lower-bound}
we have $O(\log s)^{h-1}\ge n/\log^{2}n$ and the result follows.
\end{proof}
%
\begin{proof}[Proof of \expref{Corollary}{cor:SOT-lower-bounds}] We show that these
circuits have small probabilistic degree. We approximate each And
and Or gate by a degree-$\log^{O(1)}(s)$ polynomial with error $0.1/s$
by \expref{Lemma}{lem:approx-KS}. We approximate each Sym and $\Ltf$ gate
with fan-in $m$ by a degree $\sqrt{m}\cdot\log^{O(1)}(s)$ polynomial
with error $0.1/s$ by \expref{Lemmas}{lem:AW-maj} \expref{and}{lem:Sri-LTF}.
We compose all these polynomials. By a union bound, the error is $<1/3$.

It remains to bound the degree. We prove by induction that for every
gate $g$ at distance $h'$ from the input, the composed polynomial
(obtained by composing the polynomial for each gate) has degree $\sqrt{t'}\cdot\log^{O(h')}s$,
where $t'$ is the least value such that the sub-circuit obtained
by considering $g$ as output is $t'$-SoT-$\AC^{0}[\oplus]$. This
is obvious for $h'=1$ by the above approximations. For larger $h'$,
suppose $g$ is an And or Or gate. Then by induction the sub-circuits
have degree $\sqrt{t'}\cdot\log^{O(h'-1)}s$ and the result follows.
Otherwise, $g$ is either a Sym or an $\Ltf$ gate with fan-in $m$. Hence
some child of $g$ is $t'/m$-SoT-$\AC^{0}[\oplus]$, and others are
$t''$-SoT-$\AC^{0}[\oplus]$ for $t''\le t'/m$. By induction they
all have degree $\sqrt{t'/m}\cdot\log^{O(h'-1)}s$ and so the degree
of $g$ is $\sqrt{m}\cdot\log^{O(1)}(s)\cdot\sqrt{t'/m}\cdot\log^{O(h'-1)}s$
as desired.
\end{proof}

\section{Uniform distribution: Proof of \expref{Theorem}{thm:uniform} \label{sec:Uniform-distribution:-Proof}}

In this section we prove \expref{Theorem}{thm:uniform}. We need some
tools from the theory of hardness amplification.

\subsection{Tools from hardness amplification}

First we need a relatively standard connection between average-case
and worst-case lower bounds for general circuits. The important point
here is that we are not claiming that the circuits are of constant
depth or have any other restriction except size.

%
%
\begin{lem}[{\cite{BFNW93,Imp95,STV01}}]
\label{lem:STV} Let $s=s(n)$ satisfy $s(O(n))\le s^{O(1)}(n)$.
There is a constant $c$ such that the following holds.

Suppose that for every function $f:\zo^{n}\to\zo$ in $\EEE^{\NP}$ there
is a circuit of size $s(n)\ge n$ such that $\P_{x\in\zo^{n}}[C(x)=f(x)]\ge1/2+1/n$.
Then for every function $f:\zo^{n}\to\zo$ in $\EEE^{\NP}$ there is a
circuit of size $s^{O(1)}(n)$ such that $C=f$.
\end{lem}

\begin{rem}
For the proof we can use Theorem 24 and Lemma 28 in~\cite{STV01}.
That gives the size bound $s'=n^{O(1)}\cdot s(O(n))$ which is $\le s^{O(1)}(n)$
by the assumptions on $s$. The only thing that is left to verify
is that their encoding procedure can be implemented in $\EEE^{\NP}$.
Indeed, as observed already in~\cite{DBLP:conf/latin/ChenOS18}, given
a function $f$ in $\EEE^{\NP}$ we can construct another $\EEE^{\NP}$ function
which first writes down the truth table of $f$ by calling the algorithm
for $f$ $2^{n}$ times, using the $\NP$ oracle for each call. Then we
can run the encoding procedure in~\cite{STV01} which runs in polynomial
time in the length $2^{n}$ of the truth table.
\end{rem}

We need another result from the hardness amplification literature.
This time we need a much finer bound on the complexity of the associated
decoding procedure. We need that the proof of a hardness amplification
from constant to $1/2-\e$ can be implemented essentially by a majority
on $1/\e^{2}$ bits. Such a fine result actually follows from Impagliazzo's
hard-core-set proof~\cite{Imp95} of Yao's XOR lemma (see~\cite{GNW95}),
as was made explicit in a beautiful paper by Klivans~\cite{Kli01}.
Klivans actually was working in a slightly different context and did
not need to bound the fan-in of the majority gate, but the bound is
evident in his analysis. We note that %
Klivans's   %
application (a switching-lemma
free proof of the average-case lower bounds for $\AC^{0}$) was later
simplified by Klivans and Vadhan (see the exposition in~\cite{viola-FTTCS09}).
This simplification bypasses the XOR lemma, using instead the random
self-reducibility of parity. In our application, the function may
not be randomly self-reducible.
\begin{lem}
\label{lem:Klivans}Let $f:\zo^{n}\to\zo$ be a function that is $\Omega(1)$-hard
for $\Maj_{t}\circ \AC^{0}[\oplus]$ circuits of size $s$ and depth
$h$. Then the function $f':(\zo^{n})^{k}\to\zo$ defined as $f'(x_{1},x_{2},\ldots,x_{k}):=\oplus_{i\le k}f(x_{i})$
is $(1/2-O(1/\sqrt{t})-(1-\Omega(1))^{k})$-hard for $\AC^{0}[\oplus]$
circuits of size $s/t$ and depth $h-1$.
\end{lem}

\begin{rem}
The combination of Theorem 8 and Fact 9 in~\cite{Kli01} shows that
there exists a ``hard-core'' set $S\subseteq\zo^{n}$ of size $\Omega(2^{n})$
such that for every $\AC^{0}[\oplus]$ circuit $C$ of size $s/t$
and depth $h-1$ one has $\P_{x\in S}[C(x)\ne h(x)]\ge1/2-O(1/\sqrt{t})$.
Otherwise, one can take the majority of $t$ such circuits to obtain
a contradiction. This suffices for a proof of the XOR lemma (see Lemma
4 in~\cite{Imp95} or Theorem 10 in~\cite{Kli01}).
\end{rem}


\subsection{A result by Alman and Chen}

We need the following result which is essentially Theorem 1.9 by Alman
and Chen~\cite{AlmanC19}.
\begin{lem}
\label{lem:stronger-AlmanChen-Conditional-1} Suppose that $\EEE^{\NP}$
has circuits of size $s$. Then there is a function in $\EEE^{\NP}$ that
is $\Omega(1)$-hard for $\F_{2}$-polynomials of degree $\Omega(n/\log s)$,
infinitely often.
\end{lem}

Theorem 1.9 by Alman and Chen~\cite{AlmanC19} claims $\Omega(1)$-hardness
for rank $n^{1-\alpha}$ for any $\alpha>0$. As mentioned earlier
and in~\cite{ServedioV-sparse}, this implies the same hardness for
degree \emph{$n^{1-\alpha}$} for any $\alpha>0$. We note here slightly
stronger parameters, and we give a streamlined proof for polynomials,
following closely theirs, in \expref{Section}{sec:Proof-of-Lemma-stronger-AC}.

\subsection{Proof of \expref{Theorem}{thm:uniform}}

Towards a contradiction, suppose that for every $f:\zo^{n}\to\zo$
in $\EEE^{\NP}$ and for all sufficiently large input lengths $n$ there
are $\AC^{0}[\oplus]$ circuits $C$ of size $s=s(n)$ and depth $h$
such that $\P_{x\in\zo^{n}}[f(x)\ne C(x)]\le1/2-(\log^{9h}s)/n$.

By \expref{Lemma}{lem:STV}, $\EEE^{\NP}$ has circuits of size $s^{O(1)}(n)$.

By \expref{Lemma}{lem:stronger-AlmanChen-Conditional-1}, there is a function
in $\EEE^{\NP}$ that is $\Omega(1)$-hard under the uniform distribution
for polynomials of degree $\Omega(n/\log s)$, for infinitely many
$n$. In particular, for those $n$, the latter function is $\Omega(1)$-hard
under the uniform distribution for $\MAJ_{n^{2}/\log^{8h}s}\circ \AC^{0}[\oplus]$
circuits where the $\AC^{0}[\oplus]$ circuits have size $2^{\log^{2}s(n)}$
and depth $h+1$. This is because the latter circuits have $o(1)$-error
degree $d:=O(\sqrt{n^{2}/\log^{8h}s}\cdot\ensuremath{\log^{2h}s\cdot\log n)\le o(n/\log s)}$.
This is proved by applying \expref{Lemma}{lem:AW-maj} to $\MAJ$ \expref{and}{lem:approx-KS}
with error $1/n^{2}$ to the $\AC^{0}[\oplus]$ circuits. So if such
a circuit computed $h$ correctly on a $\delta$ fraction of the inputs
there would be a distribution $P$ on polynomials of degree $d$ such
that $\P_{x,P}[P(x)=h(x)]\ge\delta-o(1)$, where $x$ is uniformly
distributed. And so by an averaging argument we could fix a specific
polynomial to contradict the claim above.

Now we can apply \expref{Lemma}{lem:Klivans} to construct a function
on $nk$ bits for $k=O(\log n)$ which, for infinitely many $n$,
has hardness %
$1/2-O(1)\Big/\sqrt{n^{2}/\log^{8h}s(n)}-1/n\ge1/2-\log^{2h}s(n)/n$ %
for circuits of size $2^{\log^{2}s(n)}/(n^{2}/\log^{O(h)}s(n))\ge2^{\Omega(\log^{2}s(n))}$,
where the inequality holds without loss of generality, and depth $(h+1)-1=h$.
Here we are using that $(1-\Omega(1))^{k}\le1/n$ for $k=O(\log n)$.

To contradict our initial assumption, it suffices that $s(O(n\log n))<2^{\Omega(\log^{2}s(n))}$.
This is indeed our assumption.

\section{Proof of \expref{Lemma}{lem:stronger-AlmanChen-Conditional-1} \label{sec:Proof-of-Lemma-stronger-AC}}

We need the following recent construction of \emph{smooth} PCPs of
proximity by Paradise~\cite{DBLP:journals/eccc/Paradise19}, which
we state using notation similar to that in \expref{Section}{subsec:PCP-tools}.
\begin{lem}
\label{lem:smooth-PCP}For any constants $\delta,\sigma\in(0,1/3]$
there is a constant $c$ and an algorithm $M$ with the following
properties. On input a circuit $C$ of size $s$ on $\ell\le s$ inputs,
$\ell$ a power of $2$, the algorithm $M(C)$ computes in time $s^{c}$
circuits $Q:\zo^{r}\to(\zo^{r})^{c}$, $Q':\zo^{r}\to(\zo^{\log\ell})^{c}$,
and $R:\zo^{2c}\to\zo$ such that:

\textbf{Proof length.} $2^{r}\le s^{c}$,

\textbf{Completeness.} If $y$ makes $C(y)=1$ then there exists a
map $\pi:\zo^{r}\to\zo$ such that for any $z\in\zo^{r}$ we have
$R(y(q'_{1}),\ldots,y(q'_{c}),\pi(q_{1}),\ldots,\pi(q_{c}))=1$ where
$(q_{1},\ldots,q_{c})=Q(z)$, $(q'_{1},\ldots,q'_{c})=Q'(z)$, and
$y(i)$ is the $i$ bit of $y$,

\textbf{Soundness.} If $y$ is $\delta$-far in Hamming distance from
any $z$ such that $C(z)=1$, then for every map $\pi:\zo^{r}\to\zo$,
at most a $\sigma$ fraction of the
$z\in\zo^{r}$ have  %
%
\[ R(y(q'_{1}),\ldots,y(q'_{c}),\pi(q_{1}),\ldots,\pi(q_{c}))=1 \]  %
where $(q_{1},\ldots,q_{c})=Q(z)$, $(q'_{1},\ldots,q'_{c})=Q'(z)$,

\textbf{Smoothness.} For every $i\le c$, each query $q_{i}$ is uniformly
distributed in $\zo^{r}$ over the choice of $z\in\zo^{r}$.
\end{lem}

\begin{proof}[Proof of \expref{Lemma}{lem:stronger-AlmanChen-Conditional-1}]
As in the
proof of \expref{Theorem}{thm:prob-degree-lower-bound},
let $L\in\mathrm{\NTime}(2^{n})\setminus\mathrm{\NTime}(o(2^{n}))$
\cite{Cook73,SeiferasFM78,Zak83} be a unary language. We want to
show that $L\in\mathrm{\NTime}(o(2^{n}))$ to reach a contradiction
(assuming the opposite of the conclusion, which we will do shortly).
Consider the algorithm $f(x,y)$ that on input $x\in\zo^{n}$ and
$y\in\{1,2,\ldots,2^{n}\poly(n)\}$ constructs the circuits from the
PCP in \expref{Lemma}{lem:BV-PCP}, computes the first satisfying assignment
if one exists, and outputs its bit $y$. This algorithm can be implemented
in $\EEE^{\NP}$ by computing the satisfying assignment one bit at the
time. By the assumption, this algorithm can be implemented by a circuit
$C$ of size $s\ge n$.

The algorithm to show %
$L\in \mathrm{\NTime}(o(2^{n}))$   %
proceeds by guessing
this circuit $C$, and composing it with the PCP in \expref{Lemma}{lem:BV-PCP}
to obtain another circuit $C'$ of size $s^{O(1)}$. By the properties
of the PCP, it then suffices to efficiently compute
\[
\P_{z\in\zo^{r}}[C'(z)=1]\,.
\]

We use the PCP in \expref{Lemma}{lem:smooth-PCP} to aid in this. We first
have to encode $z$. Fix any efficient, linear, binary error-correcting
code $ECC:\zo^{r}\to\zo^{r_{\delta}}$ with Hamming distance $\delta$
and efficient encoding and decoding, where $r_{\delta}$ is linear
in $r$ for any fixed $\delta$. By composing the decoding procedure
with $C'$ we obtain another circuit $C''$ of size $s^{O(1)}$ and
see that it suffices to compute
\[
\P_{z\in\zo^{r}}[C''(ECC(z))=1]\,.
\]

Applying the PCP in \expref{Lemma}{lem:smooth-PCP} to circuit $C''$
we see that if $x\in L$ then there are proofs $\pi_{z}$ such that
\[
\P_{z\in\zo^{r},z'\in\zo^{O(\log s)}}[R(ECC(z)(q'_{1}),\ldots,ECC(z)(q'_{c}),\pi_{z}(q_{1}),\ldots,\pi_{z}(q_{c}))=1]=1\,,
\]
while if $x\not\in L$ then for any proofs $\pi_{z}$ we have
\[
\P_{z\in\zo^{r},z'\in\zo^{O(\log s)}}[R(ECC(z)(q'_{1}),\ldots,ECC(z)(q'_{c}),\pi_{z}(q_{1}),\ldots,\pi_{z}(q_{c}))=1]\le1/2\,,
\]

where $(q_{1},\ldots,q_{c})=Q(z')$, $(q'_{1},\ldots,q'_{c})=Q'(z')$,
and picking the soundness in the PCPs in \expref{Lemmas}{lem:BV-PCP}
\expref{and}{lem:smooth-PCP} to be a small enough constant.

The proofs $\pi_{z}(q)$ can be computed as a function of $z$ and
$q$ by an $\EEE^{\NP}$ algorithm $f(z,q)$. Now let us assume the opposite
of the conclusion of the lemma. Then for every $\beta>0$ there are
polynomials $p$ of degree $\beta d$ such that $\P_{z,q}[p(z,q)\ne f(z,q)]\le\beta$.
Because the queries $q$ are smooth, and by a union bound, we have
that if $x\in L$ then
\[
\P_{z\in\zo^{r},z'\in\zo^{O(\log s)}}[R(ECC(z)(q'_{1}),\ldots,ECC(z)(q'_{c}),p(z,q_{1}),\ldots,p(z,q_{c}))=1]\ge1-\beta c\,,
\]
while if $x\not\in L$ then the above probability is $\le1/2$ as
above. Thus picking $\beta<c/2$ there is a gap in the above probabilities.
If we can detect this gap efficiently, we are done. (Recall that $L$
is unary, so membership of $x$ in $L$ is solely determined by its
length.)

The algorithm will guess the polynomial $p$ and then compute the
above probability. To compute the probability, it enumerates over
all $s^{O(1)}$ choices for $z'$. For any fixed $z'$, the function
\[
R(ECC(z)(q'_{1}),\ldots,ECC(z)(q'_{c}),p(z,q_{1}),\ldots,p(z,q_{c}))
\]
is a polynomial of degree $2c\beta d$. Here we simply write $R$
as a polynomial of degree $2c$, and use that $ECC$ is linear. Applying
\expref{Theorem}{thm:=000023SAT-poly} allows us to compute this probability
in time $2^{n-\Omega(n/\beta d)+O(\log n)}$. If $d=n/\log s$ the
running time is $2^{n-\Omega(\log s)/\beta+O(\log n)}$. Picking $\beta$
small enough the saving in time can compensate for the $s^{O(1)}$
slow-down that comes from enumerating over all $z'$, and the proof
is concluded.
\end{proof}

\section{Number-on-forehead \label{sec:Number-on-forehead}}

The connection between probabilistic degree and number-on-forehead
lower bounds was pointed out in~\cite{Viola-map}. We present next
a stronger quantitative relationship.
\begin{thm}
\label{thm:NOF-connection}Suppose that $f:\zo^{n}\to\zo$ has $(1/2-\e)$-error
degree $d$. Then $f$ has number-on-forehead protocols with $O(d\log(n/\e))$
parties and communication $O(d\log(n/\e)\cdot\log n)$, under any
partition of the input.
\end{thm}

To illustrate, if $d=\log n$ and $\e=1/\sqrt{n}$ (something which
we cannot rule out for any function in $\NP$) then we obtain $O(\log^{2}n)$
parties and $O(\log^{3}n)$ communication.
\begin{proof}
Let $D$ be the distribution on degree-$d$ polynomials witnessing
the $(1/2-\e)$-error degree of $f$. By taking the majority of $t=O(n/\e^{2})$
independent copies of $D$, we can drive the error to $<2^{-n}$.
By the probabilistic method we can fix the values of the $D$ and
obtain that there are $t$ polynomials $p_{1},\ldots,p_{t}$ of degree
$d$ such that $f(x)=\Maj(p_{1}(x),\ldots,p_{t}(x))$ for every input
$x\in\zo^{n}$.

Now view the $p_{i}$ as integer polynomials, and compose each with
a modulus amplifying polynomial $F$ of degree $O(\log t)$ given
by \expref{Lemma}{lem:=00005BModulus-amplifying-polynomials=00005D}.
This guarantees that the value of each polynomial modulo 2 is the
same as the value modulo $2^{O(\log t)}>t$. Hence, the value of
$\Maj(p_{1}(x),\ldots,p_{t}(x))$
is determined by
\[
q(x):=\sum_{i\le t}F(p_{i}(x))\,.
\]

Each $F(p_{i}(x))$ is a polynomial of degree $D=O(d\log t)$ with
$O(\log t)$-bit coefficients. Hence the same is true about $q(x)$.

By %
H{\aa}stad  %
and Goldmann's simulation \cite{HaG91}, $q(x)$ can be
computed by a number-on-forehead protocol with $D+1$ parties, under
any partition of the input with $D+1$ sets. Recall in the protocol
each party just sends the sum of the monomials in $q$ assigned to
them. Monomials are arbitrarily assigned to parties who know their
values. (For every monomial there exists at least one party who knows
all the bits in the monomial, since each party misses one of $D+1$
sets and the monomial has degree $D$.) The sum sent by a party has
magnitude at most the magnitude of one coefficient times the number
of monomials in $q$, which is $\le\poly(t)\cdot n^{O(d\log t)}$.
The number of bits is then $O(d\log(n/\e)\cdot\log n)$.
\end{proof}

\section{Data structures \label{sec:Data-structures}}

In this section we show that sufficiently strong lower bounds on probabilistic
rank imply new data-structure lower bounds. We say that a function
$f:\zo^{n}\times\zo^{\mu}\to\zo$ has a space-$s$ $t$-local data
structure if there exists a function $E:\zo^{n}\to\zo^{s}$ and $m=2^{\mu}$
functions $g_{q}$ such that for every $x,q$ we have
\[
f(x,q)=g_{q}(E(x))\,,
\]

where the functions $g_{q}$ are $t$-local, that is they depend on
only $t$ bits of their input at fixed locations. Here $x$ is the
input data of length $n$, $s$ is the \emph{space} of the data structure,
$q$ is the query, and $t$ is the time of the data structure. The
data structure is supposed to encode $x$ via $E(x)$ so that queries
can be answered fast. Such problems are known as \emph{static }(because
the data does not change over time) and \emph{non-adaptive} (because
$g_{q}$ reads bits at fixed locations).

Despite intense research, the state of known lower bounds can be summarized
with the following inequality

\begin{equation}
t\ge\log(m/n)/\log(s/n)\,.\label{eq:dslb}
\end{equation}

Specifically, no explicit function for which a bound better than
\expeqref{Equation}{eq:dslb}  %
is known, for any setting of parameters, and even for functions in
$\EEE^{\NP}$. We refer to~\cite{viola-datastructurelb-implies-cclb}
for additional discussion. We raise the question of improving this
state of affairs using the ``guess-and-sat'' method.

We note that several connections between data-structure lower bounds
and other challenges in computational complexity appear in the literature:
\cite{MiltersenNSW98,viola-datastructurelb-implies-cclb,DBLP:conf/stoc/DvirGW19,DBLP:journals/eccc/Corrigan-GibbsK18,DBLP:conf/innovations/RamamoorthyR20}.
Some of these works explain the failure to prove stronger lower bounds
by showing that such lower bounds would also solve other long-standing
problems in complexity theory. However in some cases the latter lower
bounds were actually found. For some settings of parameters, the connection
in~\cite{MiltersenNSW98} has been subsumed by the lower bounds in
\cite{Ajtai05} and~\cite{BSSV03}, and the connection in~\cite{DBLP:conf/stoc/DvirGW19}
by those in \cite{AlmanC19,DBLP:journals/corr/abs-2005-03123}.

Here we consider the reverse direction and show that data-structure
lower bounds follow from sufficiently strong probabilistic-rank lower
bounds. Specifically, it suffices to show that $f$ cannot be approximated
with error about $2^{-t}$ by matrixes with rank $s$. And in fact
it suffices to consider matrices $M=AB$ where $A$ is fixed and $B$
is sparse.
\begin{claim}
Suppose a function $f:\zo^{n}\times\zo^{\mu}\to\zo$ has a space-$s$
$t$-local data structure. Then there exists a $2^{n}\times(s+1)$
matrix $A$ and a distribution on $(s+1)\times2^{\mu}$ matrices $B$
such that any column of $B$ has at most $t+1$ non-zero elements,
and for every $(x,q)$ we have
\[
\P_{B}[(AB)_{x,q}=f(x,q)]\ge1/2+\Omega(\sqrt{2^{-t}})\,,
\]
where the matrices are considered over $\F_{2}$.
\end{claim}

To illustrate, let us consider the setting $m=n^{10}$, so $\mu=10\log n$
and $s=n^{1.1}$. In this case the bound in \expeqref{Equation}{eq:dslb}
becomes
\[
t\ge\log(n^{9})/\log(n^{0.1})
\]
which is constant. To improve this bound it would be enough to bound
above the probability in the claim above by $1/2+o(1)$, however the
corresponding matrix is highly unbalanced. At the other end of the
spectrum, consider the case $m=2^{n}$, so $\mu=n$, and $s=2n$.
In this case the bound in \expeqref{Equation}{eq:dslb} becomes
\[
t\ge\Omega(\log2^{n})=\Omega(n)\,.
\]
To improve on this we would need to bound the probability by a quantity
exponentially close to $1/2$.
\begin{proof}
We rely on the following basic fact. Let $h:\zo^{t}\to\{0,1\}$ be
a function on $t$ bits, and let $g(x):=(-1)^{h(x)}$. Then there
is a joint distribution $(S,b)$ where $S$ is a subset of $\{1,2,\ldots,n\}$
and $b\in\{0,1\}$ such that for every $x$ the expectation $\E_{S,b}[(-1)^{b+\sum_{i\in S}x_{i}}]$
is $g(x)/\sum_{\a}|\hat{g}_{\a}|$ where $g(x)=\sum_{\alpha}\hat{g}(\alpha)(-1)^{b+\sum_{i\in\a}x_{i}}$
is the Fourier expansion of $g$ (see~\cite{ODonnell14} for background).
To verify this, define the distribution so that with probability $|\hat{g}_{\a}|/\sum_{\a}|\hat{g}_{\a}|$
it yields outcome $(-1)^{b}=\sign\hat{g}_{\a})$ and $S=\a$. Then
the expectation is as claimed, since
\[
\E_{S,b}[(-1)^{b+\sum_{i\in S}x_{i}}]=
%
%
\sum_{\a}|\hat{g}_{\a}|\sign\hat{g}_{\a})(-1)^{\sum_{i\in\a}x_{i}}\Big/\sum_{\b}|\hat{g}_{\b}| %
\]
and $|\hat{g}_{\a}|\sign\hat{g}_{\a})=\hat{g}_{\a}$.

Also note $\sum_{\b}|\hat{g}_{\b}|\le\sqrt{2^{t}}\sum_{S}|\hat{g}_{\b}|^{2}=\sqrt{2^{t}}$
by Cauchy--Schwarz and Parseval, using that $g$ has range $\{-1,1\}$.
From this we deduce that $\E_{(S,b)}[(-1)^{b+\sum_{i\in S}x_{i}}g(x)]\ge1/\sqrt{2^{t}}$.
And so $\P[b+\sum_{i\in S}x_{i}=h(x)\mod2]\ge1/2+0.5/\sqrt{2^{t}}$.

Row $x$ of the matrix $A$ is $E(x)$ with an extra $1$ appended
at the end. Column $q$ of $B$ is obtained by applying the fact above
to the function $g_{q}$, setting to $1$ the bits corresponding to
$S$ and to $0$ the others, possibly except the last which is used
to implement the bit $b$.
\end{proof}
With this approach, the ``$+1$'' is unavoidable: the data structure
may choose to encode $x$ with the all-zero string, so we can never
get a non-zero value if our row is just $E(x)$.

%
\subsubsection*{Acknowledgments}   %

We thank the anonymous referees for the detailed and helpful feedback.

\bibliographystyle{tocplain}
%
%
\bibliography{bibstrings,v022a007,bibtail}  %

\begin{tocauthors}
\begin{tocinfo}[viola]
 Emanuele Viola\\
 Associate professor \\
 Khoury College of Computer Sciences\\
 Northeastern University\\
 Boston, MA\\
 viola\tocat{}ccs\tocdot{}neu\tocdot{}edu \\
 \url{http://www.ccs.neu.edu/~viola}
\end{tocinfo}

\end{tocauthors}

\begin{tocaboutauthors}
%
\begin{tocabout}[viola]  %
{\sc Emanuele Viola} 
is happy that
%
\href{https://theoryofcomputing.org}{\emph{Theory of Computing}} exists.
\end{tocabout}
\end{tocaboutauthors}

\end{document}
